% Ad Familiares 3.4 --- headnote
% ad-familiares-03-04, around 5 June 51 BC, Brundisium

\textit{Cicero to Appius Claudius Pulcher, written from
Brundisium on or about the 4th of June 51 BC (the
manuscript dateline: \textit{Scr. Brundisi prid. Non.
aut ipsis Non. Iun. a. 703}). Cicero is about to cross
the Adriatic on his way out to Cilicia; Appius is the
outgoing proconsul whom he will succeed there. The
exchange is awkward in the way that handovers between
high-magnates often were: Appius has been an active
governor of the kind Cicero himself disapproves of ---
extortionate by report, slow to vacate, sensitive about
his standing --- and the two men, college-fellows in
the augurate and political associates of long standing,
must now manage a transition that will inevitably
expose the differences in their administration.}

\textit{The tone of the letter is studiously gracious.
Cicero acknowledges three named messengers of Appius's
goodwill, picks out Appius's recent gift of a book on
augural law as a particular pleasure, and triangulates
the relationship through their shared family
connections (Pompey, father-in-law of Appius's daughter;
Brutus, Appius's son-in-law). The protestations are
warmer than the underlying business: Cicero is, in
effect, refusing to commit to anything definite until he
has heard, through L. Clodius, what Appius actually
wants, while keeping the public courtesies polished. The
correspondence with Appius, gathered in book 3 of the
\textit{Ad Familiares}, will run on through the year in
much this register --- the polite cooling of a
relationship that survives but no longer fully trusts.}
